\section{What projection measures and steering changes}
\label{sec:calibration}

\subsection{Register dominates the readout}

On Gemma 3 4B, the axis contains an identity signal, but writing register
has a larger effect on the projections. We create four conditions by
combining human or AI identity with warm or formal writing. Across these
conditions, register shifts the layer-33 projection about three times as far
as identity ($+337$ vs. $+114$; Table~\ref{tab:register2x2} in
Appendix~\ref{app:register2x2factorial}). A formal response with no identity
claim reaches $+299$ points, farther AI-ward than the formal AI condition
($+265$). The extreme instructions partly break character in the human
conditions, so the 3:1 register-to-identity ratio is an upper bound.

Register nevertheless defines a different direction. Removing the register
direction leaves 87\% of the axis norm at both readout layers, and
identity-classification AUC barely drops at layer 33, from 92.6\% to 91.1\%
(Appendix~\ref{app:register2x2}). A sparse-autoencoder readout
\citep{bricken2023monosemanticity, templeton2024scaling} shows partial
feature-level overlap with register at layers 27 and 29; the layer-33
dictionary is unusable (Appendix~\ref{app:sae}). Register directions can themselves steer
\citep{konen2024style}, so we rely on matched-register comparisons and
behavioral intervention, rather than projection magnitude alone, to
interpret the axis.

\subsection{The steering effect is not reducible to register}

A register-confounded projection does not imply that axis steering is
reducible to register. At the steering layer, the axis is nearly collinear
with register (cosine 0.980) and PC1 ($\sim$0.98,
Appendix~\ref{app:register2x2decomp}). The layer-13 cosine should be read
with caution (Section~\ref{sec:axis}). Steering with a matched-norm register
direction also crosses the identity midpoint on the 10 held-out questions,
so geometry alone cannot determine what steering changes
(Appendix~\ref{app:register2x2}). To see whether steering's effect comes
from register or from the rest of the axis, we decompose it into components
parallel and orthogonal to a given direction \citep{arditi2024refusal},
doing this separately for PC1 and for register, and test all four resulting
components on prompt-injection execution (Section~\ref{sec:injection}) and
agentic prose length.

The register-parallel component contains 98\% of the axis norm at layer 13,
yet prompt-injection execution remains at baseline across doses, a measured
null. By contrast, the much smaller register-orthogonal component increases
execution with dose (supporting but not decisive evidence); a
norm-matched full register direction also remains at baseline. Within this
tested register-direction evaluation, the behavioral effect of axis steering
is therefore not explained by writing register (Table~\ref{tab:arms} and
Figure~\ref{fig:calibration}b, Appendix~\ref{app:register2x2decomp}).

PC1 gives a less clean decomposition: the PC1-parallel component has a
measurable effect, and at matched median prose the PC1-orthogonal component can produce more
execution than the full axis, though this comparison is post hoc and depends on how PC1 is estimated
(Figure~\ref{fig:calibration}c, Appendix~\ref{app:register2x2decomp}). PC1
therefore carries part, but not all, of the steering effect. The refusal-overlap
mechanism proposed by \citet{li2026pitfalls} fails our single-vector raw and
PC1-separated geometric tests; no causal ablation was run (Appendix~\ref{app:refusal}).
Within this one-model, one-layer specificity check, identity remains the best-supported
interpretation of what steering changes.
